\section*{Chapter \ref{ch:pl:data-quality-and-lineage} --- Data Quality and Lineage}

\begin{solution}{exo:pl:data-quality-and-lineage:1}
Four: 91\,687 to 91\,690. The capture of chapter~2 waits for the other line to fill them, then records the gap's range and moves
on; the sequence rule of this chapter flags the first row after it, so that consumers know the day is incomplete there.
\end{solution}

\begin{solution}{exo:pl:data-quality-and-lineage:2}
On a large-tick instrument most consecutive trades print at the same price, so the median absolute deviation of recent deviations is
zero, and any move, even one tick, would be infinitely many robust units away: every tick change would be a spike. The floor (0.1\%
here) says that nothing smaller than a plausible move is ever an anomaly.
\end{solution}

\begin{solution}{exo:pl:data-quality-and-lineage:3}
A decimal slip: a factor of exactly ten. The forward median is still 300\,000, because the next trades are back at the level, so the
print deviates from both medians; a genuine move to 3\,000\,000 would stay there and deviate only from the backward one.
\end{solution}

\begin{solution}{exo:pl:data-quality-and-lineage:4}
A flag can be wrong (the jump), a row that is wrong for one use is needed by another, and a deletion can be neither audited nor
undone. Surveillance and trade-reconstruction systems need the erroneous and broken trades: they are part of what happened.
\end{solution}

\begin{solution}{exo:pl:data-quality-and-lineage:5}
Kept: the decimal slips that multiply a price by ten dominate the volume-weighted average (a few trades at ten times the price). With
error rows excluded, the remaining $-0.02\%$ comes from the trades lost in the gap of 200 events: the sample is complete except there,
and no flag can restore missing data.
\end{solution}

\begin{solution}{exo:pl:data-quality-and-lineage:6}
The versions built from the old close: the marks, the P\&L and the research features, which used the close. The minute bars, built from
trades only, need not be rebuilt. The lineage knows because each version records the content hashes of its inputs: everything
downstream of the old close's version is stale.
\end{solution}

\begin{solution}{exo:pl:data-quality-and-lineage:7}
Naive: 24 false spike flags (the trades after SIM2's 3\% jump, until the backward median catches up) and one false staleness flag
(SIM2's 60-second halt). Aware: none; every planted defect is found by both.
\end{solution}

\begin{solution}{exo:pl:data-quality-and-lineage:8}
It deletes genuine moves (a 5\% jump is news on a volatile day, and every trade after it until the next trade) and invents quotes
the market never showed (a halted instrument did not quote; repeating the last quote makes it look tradeable). The raw data are gone,
so neither error can be found later. Flag instead, with rules that know the market's state and confirm spikes both ways, and keep
the data as delivered.
\end{solution}

\begin{solution}{pb:pl:data-quality-and-lineage:1}
\begin{enumerate}
\item Completeness (sequence gaps, stale quotes), validity (positive prices and sizes), consistency (crossed book, price spikes, the
vendor's trade count), timeliness (the contract's delivery times; the staleness rule is its intraday cousin).
\item The producer's and consumers' agreement on schema, rules and delivery times: a change becomes a negotiated change, not a
surprise.
\item 47\,472 events, 9\,258 quotes and 1\,190 trades.
\item The broken trades: they are valid prices on the day; only the venue's list of cancellations, delivered later, identifies them.
\item It is not a failure of the day's data but a new version of a fact: it is stored as a version, and the lineage finds what to
rebuild.
\item The backward median lags the new level, so each new trade deviates by about 3\% in robust units of a few basis points: 24 flags
on seed~0.
\item 495 naive, 0 aware.
\item Sixty seconds without a quote exceeds 30 seconds; the aware rule does not count the halted time.
\item Deleted two dozen genuine trades after the news, and understated the day's move.
\item A rule that finds every planted defect may still flag hundreds of genuine rows; only genuine look-alikes measure false flags.
\item Kept: SIM1 $+2.69\%$, SIM2 $+1.08\%$; excluded: within 0.03\% (SIM2 $-0.02\%$).
\item About 24\,000.
\item 16 flagged minutes, those where the gap lost trades; it says where to look, not which trades are missing.
\item A flag marks rows and leaves the partition readable; a \omterm{def:pl:data-quality-and-lineage:flag}{quarantine} withholds the partition until a person releases it, with a
reason, or it is rebuilt.
\item The marks, the P\&L and the features.
\item \textbf{Named result.} Over twenty seeds both suites find all 340 planted defects; the naive suite also flags 495 genuine trades
as spikes and 20 halts as stale feeds, the aware one none. With the defects kept, the VWAPs are 2.69\% and 1.08\% too high and SIM2's
realised variance 24\,000 times too large; with error rows excluded, the VWAPs are within 0.03\% and the variance is exact.
\item Confirming spikes against the forward median: 495 of the 515 false flags.
\item Without the rule, 20\,000 trades 60\% away from their prices, for years; with it, the same trades flagged as broken, excluded from
research and kept for surveillance.
\item For each dataset version, the content hashes of the input versions it was built from (and the code), so that a new input
version names everything downstream.
\item A data-quality layer records what it found and lets people decide; a cleaning script changes the data and leaves no trace.
\end{enumerate}
\end{solution}

\begin{solution}{iq:pl:data-quality-and-lineage:1}
Sequence gaps, validity bounds, crossed or locked quotes, spikes confirmed both ways, stale instruments while the market is open,
trades against the venue's list of cancellations, counts and volumes against a second source and the official close; each with an
owner, and results as flags.
\iqlookfor{The four families, the market's state, flags not deletions.}
\end{solution}

\begin{solution}{iq:pl:data-quality-and-lineage:2}
A bad print usually comes back: compare it with the prices before and after, in robust units with a floor; a genuine move persists.
Check also against quotes, other venues and the news, and against the venue's cancellations published later.
\iqlookfor{Look-ahead in a batch check; reversal as the signature.}
\end{solution}

\begin{solution}{iq:pl:data-quality-and-lineage:3}
The rule may be wrong, other consumers need the rows, and deletion cannot be audited or reversed. Flags let each consumer choose and
let the rules improve.
\iqlookfor{Reversibility and different consumers.}
\end{solution}

\begin{solution}{iq:pl:data-quality-and-lineage:4}
Store the correction as a new version of the close, then use the lineage to find every result built from the old version (marks, P\&L,
risk, features) and rebuild them, with the P\&L change reported as a correction, not silently replaced.
\iqlookfor{Versions and lineage; the change reported.}
\end{solution}

\begin{solution}{iq:pl:data-quality-and-lineage:5}
Plant known defects and, just as important, genuine look-alike events (halts, jumps, auctions, early closes), and measure detection
and false flags per rule; track flag rates over time on real data and review a sample of flags each week.
\iqlookfor{Confounders and false-flag rates, not only detection.}
\end{solution}

\begin{solution}{iq:pl:data-quality-and-lineage:6}
Name every dataset version by content hash; record for each its input versions, code version, parameters and run; store results with
the manifest; answer ``what is this built from'' and ``what depends on this'' by graph queries; recompute from the manifest after a
correction.
\iqlookfor{Content-addressed versions and a queryable graph.}
\end{solution}
